\documentclass[9pt,a4paper,twocolumn]{ctexart}

% --- Packages ---
\usepackage[top=1.2cm, bottom=1.2cm, left=1.2cm, right=1.2cm, columnsep=0.8cm]{geometry}
\usepackage{hyperref}
\usepackage{graphicx}
\usepackage{float}
\usepackage{longtable}
\usepackage{tabularx}
\usepackage{booktabs}
\usepackage{enumitem}
\usepackage{xcolor}
\usepackage{fancyhdr}
\usepackage{listings}
\usepackage{fontspec}
\usepackage{titlesec}
\usepackage{caption}

% --- Code block styling ---
\definecolor{codebg}{RGB}{245,245,245}
\definecolor{codeframe}{RGB}{220,220,220}
\definecolor{codekw}{RGB}{0,0,192}
\definecolor{codecomment}{RGB}{0,128,0}
\definecolor{codestring}{RGB}{163,21,21}
\definecolor{codenum}{RGB}{128,0,128}
\definecolor{codepunct}{RGB}{90,90,90}
\definecolor{badgegreen}{RGB}{34,139,34}
\definecolor{badgeblue}{RGB}{30,144,255}
\definecolor{badgeblack}{RGB}{50,50,50}

% --- Difficulty badges (rounded corners via tikz, pre-rendered in saveboxes) ---
\usepackage{tikz}
\newsavebox{\badgechubox}
\savebox{\badgechubox}{\tikz[baseline]{\node[fill=badgegreen,text=white,rounded corners=2.5pt,inner xsep=4pt,inner ysep=1.5pt,font=\sffamily\footnotesize\bfseries]{初级};}}
\newsavebox{\badgezhongbox}
\savebox{\badgezhongbox}{\tikz[baseline]{\node[fill=badgeblue,text=white,rounded corners=2.5pt,inner xsep=4pt,inner ysep=1.5pt,font=\sffamily\footnotesize\bfseries]{中级};}}
\newsavebox{\badgegaobox}
\savebox{\badgegaobox}{\tikz[baseline]{\node[fill=badgeblack,text=white,rounded corners=2.5pt,inner xsep=4pt,inner ysep=1.5pt,font=\sffamily\footnotesize\bfseries]{高级};}}
\newcommand{\lvchu}{\usebox{\badgechubox}}
\newcommand{\lvzhong}{\usebox{\badgezhongbox}}
\newcommand{\lvgao}{\usebox{\badgegaobox}}
\pdfstringdefDisableCommands{%
  \def\lvchu{[初级]}%
  \def\lvzhong{[中级]}%
  \def\lvgao{[高级]}%
}

\lstset{
  basicstyle=\ttfamily\scriptsize,
  keywordstyle=\color{codekw}\bfseries,
  commentstyle=\color{codecomment}\itshape,
  stringstyle=\color{codestring},
  backgroundcolor=\color{codebg},
  frame=single,
  rulecolor=\color{codeframe},
  breaklines=true,
  breakatwhitespace=false,
  breakindent=0pt,
  postbreak=\mbox{\textcolor{red}{$\hookrightarrow$}\space},
  numbers=left,
  numberstyle=\tiny\color{gray},
  columns=flexible,
  keepspaces=true,
  showstringspaces=false,
  tabsize=2,
  aboveskip=4pt,
  belowskip=4pt,
  extendedchars=true,
  literate={，}{{，}}1 {；}{{；}}1 {！}{{！}}1 {？}{{？}}1 {“}{{“}}1 {”}{{”}}1 {（}{{（}}1 {）}{{）}}1 {《}{{《}}1 {》}{{》}}1,
}

% --- Extra language definitions (no built-in listings support) ---
\lstdefinelanguage{json}{
  morekeywords={true,false,null},
  sensitive=true,
  morestring=[b]",
  comment=[l]{//},
  morecomment=[s]{/*}{*/},
}
\lstdefinelanguage{yaml}{
  morekeywords={true,false,null,yes,no,on,off},
  sensitive=false,
  morecomment=[l]{\#},
  morestring=[b]",
  morestring=[d]',
}
\lstdefinelanguage{http}{
  morekeywords={GET,POST,PUT,DELETE,PATCH,HEAD,OPTIONS,HTTP,Host,Accept,Authorization,Content-Type,Content-Length,User-Agent},
  sensitive=true,
  morecomment=[l]{\#},
  morestring=[b]",
}
\lstdefinelanguage{TypeScript}{
  morekeywords={abstract,any,as,async,await,break,case,catch,class,const,continue,debugger,declare,default,delete,do,else,enum,export,extends,false,finally,for,from,function,get,if,implements,import,in,instanceof,interface,keyof,let,module,namespace,never,new,null,number,of,package,private,protected,public,readonly,return,set,static,string,super,switch,symbol,this,throw,true,try,type,typeof,undefined,unknown,var,void,while,with,yield,Record,Array,Promise,AsyncIterable,ReadableStream,Date,Error,Object,JSON},
  sensitive=true,
  morecomment=[l]{//},
  morecomment=[s]{/*}{*/},
  morestring=[b]",
  morestring=[b]',
  morestring=[b]`{`},
}

% --- Compact spacing ---
\linespread{0.95}
\setlength{\parskip}{1pt plus 1pt minus 1pt}
\setlength{\parindent}{0pt}

% --- Table styling ---
\renewcommand{\arraystretch}{1.1}

% --- Heading styling (numbered) ---
\titleformat{\section}{\normalsize\bfseries}{\thesection\quad}{0em}{}
\titlespacing{\section}{0pt}{6pt}{2pt}
\titleformat{\subsection}{\small\bfseries}{\thesubsection\quad}{0em}{}
\titlespacing{\subsection}{0pt}{4pt}{1pt}
\titleformat{\subsubsection}{\small\bfseries}{\thesubsubsection\quad}{0em}{}
\titlespacing{\subsubsection}{0pt}{3pt}{1pt}

% --- Page style ---
\pagestyle{fancy}
\fancyhf{}
\fancyhead[L]{\small\emph{\leftmark}}
\fancyhead[R]{\small\thepage}
\renewcommand{\headrulewidth}{0.4pt}
\renewcommand{\sectionmark}[1]{}  % prevent sections from overwriting leftmark

% --- Hyperref setup ---
\hypersetup{
  colorlinks=true,
  linkcolor=blue,
  urlcolor=blue,
  citecolor=blue,
}

% --- Caption format ---
\DeclareCaptionFormat{myformat}{\small\bfseries #1#2#3}
\captionsetup{format=myformat}

\begin{document}

\title{Agent 术语表与概念边界}
\date{}
\maketitle
\markboth{TypeScript 版 Agent 知识}{ TypeScript 版 Agent 知识}

\section*{术语表}


\subsection*{基础与系统}

{\footnotesize
\begin{tabular}{|p{\dimexpr 0.151\columnwidth-2\tabcolsep\relax}|p{\dimexpr 0.496\columnwidth-2\tabcolsep\relax}|p{\dimexpr 0.353\columnwidth-2\tabcolsep\relax}|}
\hline
\textbf{术语} & \textbf{一句话定义} & \textbf{容易混淆的点} \\
\hline
AI Agent & 能感知环境、制定计划、调用工具、执行动作并在反馈中持续运行的软件系统 & 不是所有接 LLM 的应用都是 Agent \\
\hline
Agent Loop & Agent 的核心运行循环：读取上下文 → LLM 推理 → 调用工具 → 观察结果 → 继续，直到完成或触发停止条件 & Loop 是机制，ReAct 是使用 Loop 的一种范式 \\
\hline
Chatbot & 以对话为主要交互形式、通常不主动调用外部工具改变外部世界的应用 & 有对话界面不等于有 Agent \\
\hline
Workflow & 预先定义节点与边的固定执行流程，LLM 只作为其中若干节点 & 控制权在图结构里，不在模型手里 \\
\hline
Agentic Workflow & 全局用 Workflow 管住结构，局部不确定节点嵌入 Agent 子循环的混合形态 & 不是“纯 Agent” \\
\hline
RAG & 检索增强生成：从外部知识源检索相关片段，再交给 LLM 生成回答 & RAG 是知识接入方式，不是记忆系统的全部 \\
\hline
LLM & 大语言模型，负责理解与生成 & 不是系统本身 \\
\hline
Function Calling & 模型输出结构化“调用意图”，由外部程序决定是否执行并返回结果的机制 & 模型只生成意图，不直接执行函数 \\
\hline
Tool & 让 LLM 影响外部世界的可执行能力：查库、发信、跑代码等 & Tool 是能力，Function Calling 是调用协议 \\
\hline
Tools 注册 & 把工具名、JSON Schema、用途和禁用场景描述注入上下文的工程动作 & description 质量决定调用准确率 \\
\hline
JSON Schema & 结构化描述参数类型、必填项与约束的契约 & 它是数据格式，不是通信协议 \\
\hline
Structured Outputs & 模型按给定 Schema 输出结构化结果的约束能力 & 它约束输出，不约束执行 \\
\hline
MCP & Model Context Protocol：基于 JSON-RPC 2.0 的工具/资源/提示词接入协议 & MCP 解决“怎么接入”，Function Calling 解决“怎么表达调用意图” \\
\hline
Agent Skills & 以 SKILL.md 为核心的可复用经验包，延迟加载完整指令 & Skills 是白盒经验包，Toolkit 是黑盒高阶函数 \\
\hline
Prompt Engineering & 通过改写提示词直接影响模型输出的工程方法 & 只解决“单次怎么问” \\
\hline
\end{tabular}
}



\subsection*{能力与工程}

{\footnotesize
\begin{tabular}{|p{\dimexpr 0.200\columnwidth-2\tabcolsep\relax}|p{\dimexpr 0.495\columnwidth-2\tabcolsep\relax}|p{\dimexpr 0.305\columnwidth-2\tabcolsep\relax}|}
\hline
\textbf{术语} & \textbf{一句话定义} & \textbf{容易混淆的点} \\
\hline
Context Engineering & 管理模型在每轮决策时看到什么信息的系统方法，包括系统提示词、记忆、工具描述、检索内容的动态组装 & 解决“每轮该喂什么” \\
\hline
Harness Engineering & 围绕模型外部的执行环境做工程：循环、工具、权限、上下文、可观测、错误处理 & 瓶颈常在 Harness，不在模型 \\
\hline
Memory & Agent 保存和检索历史与经验的能力，分短期/长期 & Memory 是“记住了什么”，RAG 是“去外部查什么” \\
\hline
Short-Term Memory & 当前任务范围内的对话与执行轨迹，通常放上下文 & 超出上下文即开始失真 \\
\hline
Long-Term Memory & 跨任务持久化的用户偏好、事实和经验 & 需要检索、更新、遗忘机制 \\
\hline
Planning & 把目标拆成步骤并动态调整的能力 & 计划不是写死的流程 \\
\hline
ReAct & 推理与行动交替进行的范式 & 适合路径不确定的任务 \\
\hline
Plan-and-Execute & 先全局计划、再分步执行的范式 & 适合长任务但动态调整弱 \\
\hline
Reflection & 让模型对结果进行审查和迭代改进的机制 & 不改模型权重 \\
\hline
Multi-Agent & 多个 Agent 分工协作完成任务的架构 & 通信和调试成本远高于单 Agent \\
\hline
A2A & Agent 间通信的结构化接口约定，传递带 Schema 的任务、状态与验收信息 & 不是“让 Agent 用自然语言聊天” \\
\hline
Graph & 由 Node、Edge、State 组成的工作流数据结构 & Workflow 是思想，Graph 是表达 \\
\hline
Loop & Graph 中允许回溯和重试的回边 & 必须配终止条件和成本上限 \\
\hline
LLM Gateway & 统一接入、路由、降级、限流、成本统计、审计的模型调用基础设施 & 不只是反向代理 \\
\hline
\end{tabular}
}



\section*{四组关键边界}


\subsection*{1. Agent / Chatbot / Workflow / RAG}

{\footnotesize
\begin{tabular}{|p{\dimexpr 0.091\columnwidth-2\tabcolsep\relax}|p{\dimexpr 0.227\columnwidth-2\tabcolsep\relax}|p{\dimexpr 0.182\columnwidth-2\tabcolsep\relax}|p{\dimexpr 0.295\columnwidth-2\tabcolsep\relax}|p{\dimexpr 0.205\columnwidth-2\tabcolsep\relax}|}
\hline
\textbf{维度} & \textbf{Chatbot} & \textbf{Workflow} & \textbf{Agent} & \textbf{RAG} \\
\hline
核心问题 & 怎么把话说好 & 流程是否按图执行 & 下一步该做什么 & 去哪里找相关知识 \\
\hline
控制权 & 模型直接回复 & 图结构预定义 & 模型动态决策 & 检索器 + 模型 \\
\hline
外部动作 & 通常无 & 固定节点执行 & 按需调用工具 & 只读检索为主 \\
\hline
失败模式 & 答错话 & 节点/边设计错误 & 死循环、乱调工具、目标漂移 & 召回差、上下文污染 \\
\hline
典型产品 & 客服 FAQ 机器人 & 文章审核流水线 & 故障排查 Agent & 企业知识库问答 \\
\hline
\end{tabular}
}



一句话边界：\textbf{Chatbot 重对话，Workflow 重流程，Agent 重自主决策，RAG 重知识接入}；四者可以组合，但不能互相替代。



\subsection*{2. Function Calling / MCP / HTTP API / Skills}

{\footnotesize
\begin{tabular}{|p{\dimexpr 0.348\columnwidth-2\tabcolsep\relax}|p{\dimexpr 0.413\columnwidth-2\tabcolsep\relax}|p{\dimexpr 0.239\columnwidth-2\tabcolsep\relax}|}
\hline
\textbf{层} & \textbf{解决什么} & \textbf{类比} \\
\hline
JSON Schema & 参数“长什么样” & 接口的字段定义 \\
\hline
Function Calling & 模型“想调什么、参数是什么” & 调用请求的生成协议 \\
\hline
HTTP API & 服务之间“怎么传输” & REST 接口 \\
\hline
MCP & AI 应用与工具之间“怎么发现和接入” & USB-C 统一接口 \\
\hline
Agent Skills & 一类任务“怎么做、何时做” & 团队经验包 / SOP \\
\hline
\end{tabular}
}



判断口诀：\textbf{模型生成的是调用意图，系统决定是否执行；MCP 管接入，Skills 管经验。}



\subsection*{3. Memory / RAG / Context}

{\footnotesize
\begin{tabular}{|p{\dimexpr 0.472\columnwidth-2\tabcolsep\relax}|p{\dimexpr 0.361\columnwidth-2\tabcolsep\relax}|p{\dimexpr 0.167\columnwidth-2\tabcolsep\relax}|}
\hline
\textbf{概念} & \textbf{回答的问题} & \textbf{生命周期} \\
\hline
Context & 这一轮模型能看到什么 & 单次调用 \\
\hline
Short-Term Memory & 这个任务刚才发生了什么 & 当前任务 \\
\hline
Long-Term Memory & 过去积累了什么偏好与经验 & 跨任务持久化 \\
\hline
RAG & 外部知识库里有什么相关材料 & 按需检索 \\
\hline
\end{tabular}
}



长期记忆可以借用 RAG 技术实现，但长期记忆还包括偏好、任务轨迹、经验教训的写入与演化；RAG 本身不做记忆演化。



\subsection*{4. Workflow / Graph / Loop / Agentic Workflow}

\begin{itemize}[itemsep=1pt, topsep=2pt]
  \item Workflow 是“固定流程”的抽象；
  \item Graph 是 Workflow 的数据结构表达：Node 执行、Edge 控制流、State 共享上下文；
  \item Loop 是 Graph 中的回边，用于重试、回溯、迭代；
  \item Agentic Workflow 是“固定骨架 + 局部自主”的组合。
\end{itemize}


\section*{概念关系图}


\begin{lstlisting}
flowchart LR
  User[用户] -->|任务| Agent[AI Agent]
  Agent --> LLM[LLM 推理与决策]
  Agent --> MEM[Memory 记忆]
  Agent --> TOOL[Tools 工具层]
  TOOL --> FC[Function Calling 调用意图]
  TOOL --> MCP[MCP 工具接入]
  TOOL --> SK[Agent Skills 经验包]
  Agent --> WF[Agentic Workflows]
  WF --> GRAPH[Workflow / Graph / Loop]
  MEM --> RAG[RAG 知识检索]
  AGW[Agent Gateway] --> LLM
\end{lstlisting}



\section*{FAQ}


\textbf{Q1：接了一个搜索工具的 Chatbot 是不是 Agent？} 如果每轮路径由代码写死、LLM 不决定“下一步做什么”，更接近 Workflow；只有当 LLM 能根据环境反馈持续决策并行动时，才算 Agent。



\textbf{Q2：RAG 是不是长期记忆？} 不是。RAG 是“查知识”，长期记忆还包括“记住偏好、经验、失败教训并更新”。RAG 只是长期记忆的一种实现手段。



\textbf{Q3：MCP 和 Function Calling 是不是二选一？} 不是。Function Calling 表达模型调用意图，MCP 标准化工具接入与发现。一个 MCP Server 暴露的工具，仍需用 JSON Schema 描述参数。



\textbf{Q4：有了 Agent Skills 还需要 Tool 吗？} 需要。Skills 组织经验与流程，Tool 提供原子能力；Skills 内部通常会调用 Tool。



\textbf{Q5：Multi-Agent 一定比单 Agent 强吗？} 不一定。多 Agent 带来并行与专业化，但通信、调试、成本会明显上升。先证明单 Agent 无法满足，再上多 Agent。



\section*{使用约定}

\begin{enumerate}[itemsep=1pt, topsep=2pt]
  \item 写文档时优先使用本术语表定义；出现新术语先补充到术语表再使用。
  \item 架构图、设计文档中的“Agent”不得同时混用 Chatbot 与 Workflow 含义。
  \item 涉及协议时明确写全称：MCP、A2A、JSON Schema，不自行造词。
  \item 术语变更必须同步更新 \texttt{知识地图总索引.md} 中对应节点。
\end{enumerate}


\section*{相关文档}

\begin{itemize}[itemsep=1pt, topsep=2pt]
  \item \href{../02-Agent/01-Agent核心概念.md}{AI Agent 核心概念}：Agent Loop、范式、MCP/A2A/Skills
  \item \href{../02-Agent/02-Agent记忆系统.md}{AI Agent 记忆系统}：Memory 与 RAG 的详细区别
  \item \href{../01-LLM基础/02-大模型结构化输出.md}{大模型结构化输出}：JSON Schema 与 Function Calling 落地
  \item \href{02-Agent自主性分级.md}{Agent 自主性分级}：从规则自动化到全自主的等级划分
\end{itemize}

\end{document}
